\documentclass[12pt]{letter}
\usepackage[margin=1in]{geometry}
\usepackage[T1]{fontenc}

\signature{Brent Minchew, Colin Meyer, et al.}
\address{}

\begin{document}

\begin{letter}{Dr.\ Jesse Smith\\Science\\American Association for the Advancement of Science}

\opening{Dear Dr.\ Smith,}

We submit ``Sea level is likely to rise more than 1 meter in the 21st century'' for consideration as a Research Article in \textit{Science}.

There is growing recognition that IPCC AR6 sea-level projections underestimate observed trends. We quantify the discrepancy. Using reduced physical models calibrated independently against decades of observations for each sea-level contributor---ocean thermal expansion, glaciers, Greenland, and Antarctica---we produce a component-level forecast that closes 93--101\% of observed global mean sea level over the satellite era without using GMSL as a calibration target. This is, to our knowledge, the tightest published budget closure for a component-based sea-level framework, and it serves as a genuine out-of-sample validation of the forecast system.

The forecast projects a 73\% probability of exceeding 1 meter of sea-level rise by 2100, with medians 66--72\% above IPCC medium-confidence estimates for the same warming trajectories. Our medians fall within the IPCC's own low-confidence ranges---the outcomes AR6 flags as possible but does not assign probabilities to. This paper assigns those probabilities.

A variance decomposition reveals that 92\% of total forecast uncertainty originates from a single source: the West Antarctic Ice Sheet. The remaining components---temperature-sensitive, observationally constrained, and reducible by emissions cuts---account for only 8\%. These two dimensions of risk are approximately orthogonal: emissions reduction compresses the predictable central estimate but leaves the WAIS-dominated tail largely unchanged. We convert these uncertainties from meters into people and dollars, and provide a value-of-information estimate showing that resolving the WAIS question is worth trillions of dollars in avoided costs---framing the research investment question as a cost-benefit calculation rather than an appeal. We close with a brief discussion of direct intervention on WAIS as a complement to observation and modeling, should the evidence continue to favor rapid ice loss.

The paper is approximately 4,000 words. We believe the logical chain from forecast methodology through variance decomposition to value-of-information analysis requires the full treatment in the main text, and we respectfully request consideration under \textit{Science}'s extended format.

\closing{Sincerely,}

\end{letter}
\end{document}
